%==============================================================================
%  CHAPTER 1 — INTRODUCTION  (annotated skeleton; prose to be written by author)
%------------------------------------------------------------------------------
%  HOW TO USE THIS FILE
%   - Each \section below is a STORY BEAT, not necessarily a printed heading.
%     In the final draft you may merge beats and drop some \section{} headers so
%     the Introduction reads as flowing prose. They are split here for scaffolding.
%   - "% BEAT n ..."      = the narrative job this beat does (why it's here).
%   - "% STORY:"          = the one-line thread that carries the reader forward.
%   - [WRITE: ...]        = expand into prose (this is your job; the spine is here).
%   - [FACTS: ...]        = ammunition already verified (numbers from writing.md /
%                           sources). Mark screening/interim numbers as such.
%   - \citep / \citet      = author-year placements (natbib / biblatex authoryear).
%     Keys resolve against introduction_refs.bib (merge into your references.bib;
%     remap keys if yours differ).
%
%  THE ARC IN ONE LINE
%   Clearing concentrates risk in CCPs and their members (1.1) -> clearing does not
%   destroy risk, it TRANSFORMS counterparty risk into liquidity risk through
%   procyclical margin (1.2) -> the load-bearing, under-studied node is the
%   CLIENT-CLEARING TIER (1.3) -> in stress, procyclical margin calls cascade down
%   the tier into CLIENT PAYMENT DELAYS that do not arrive independently but CLUSTER
%   / self-excite (1.4) -> studying this propagation needs a calibrated market PLUS a
%   regulation-grounded clearing tier, coupled: an ABM, and THAT MODEL IS THE
%   CONTRIBUTION (1.5) -> RQ (1.6) -> contributions (1.7) -> roadmap (1.8).
%
%  HONESTY FLAGS (writing.md conventions — keep these true in the prose)
%   * The Hawkes / payment-delay self-excitation layer is the thesis's framing lens
%     and (per writing.md) is NOT yet in the built model. Present it as the mechanism
%     the RQ foregrounds and as a stated contribution/extension — do NOT report
%     results from it unless/until implemented. See the comment in 1.4 and 1.7.
%   * Final baseline theta is PENDING a high-res run; any preview numbers are interim
%     / screening-resolution. Label them.
%   * Hill (not kurtosis) is the tail measure. Cite a reference for every modelling
%     claim; flag deviations from the Simudyne ODD explicitly.
%==============================================================================

\chapter{Introduction}
\label{ch:intro}

%==============================================================================
% BEAT 1 — HOOK: clearing was moved to the centre of the system; it mutualises
%               risk but in doing so CONCENTRATES it.
% STORY: open wide and end on the tension — the cure for bilateral counterparty
%        risk created a new concentration of it.
%==============================================================================
\section{Central clearing and the concentration of systemic risk}
\label{sec:intro-motivation}

% [WRITE: 1 short para] Post-2008 framing.
After the 2008 crisis, the G20 mandated central clearing of standardised
over-the-counter derivatives, placing central counterparties (CCPs) at the
structural centre of modern financial markets.
% [WRITE: what a CCP does, in two moves] interposition + mutualisation.
[WRITE: A CCP interposes itself between the two sides of every cleared trade
(novation), multilaterally nets exposures, and mutualises default losses through
posted margin and a shared default fund. It converts a dense web of bilateral
exposures into a hub-and-spoke structure.] \citep{menkveld_vuillemey_2021,duffie_zhu_2011}

% [WRITE: the turn — the upside is also the systemic cost] This is the tension that
% powers the whole thesis.
[WRITE: State the trade-off plainly. The same interposition that removes bilateral
counterparty risk concentrates it: the CCP becomes a node whose failure is
systemic, and the network is not unambiguously safer for every participant.]
\citep{borovkova_2013,duffie_zhu_2011}
% [FACTS: netting/insurance/fire-sale rationales and the "is a CCP safer?" debate
%  are surveyed in menkveld_vuillemey_2021; Borovkova shows safety depends on
%  network position.]

%==============================================================================
% BEAT 2 — REFRAME: clearing does not eliminate risk; it TRANSFORMS counterparty
%               risk into LIQUIDITY risk, via (procyclical) margin.
% STORY: this is the conceptual pivot of the thesis. Cont (2017) is the spine.
%==============================================================================
\section{Risk transformation: from counterparty risk to liquidity risk}
\label{sec:intro-risktransformation}

% [WRITE: the central idea] The Cont (2017) thesis, in your own words.
[WRITE: Central clearing does not make risk disappear — it transforms it. Margin
calls turn mark-to-market (accounting) losses into realised cash demands, so
counterparty risk is re-expressed as a demand on participants' liquidity buffers.]
\citep{cont_2017}

% [WRITE: why this bites in stress] procyclicality.
[WRITE: Initial margin is risk-based and therefore procyclical: it rises precisely
when volatility rises, so a stress event simultaneously raises the cash a
participant must find and drains the market in which they would raise it — the
funding/market-liquidity spiral.] \citep{brunnermeier_pedersen_2009,murphy_2014}

% [WRITE: ground it in the most recent live episode] March 2020.
[WRITE: Make it concrete with the March-2020 dash-for-cash, which is also the
thesis's stressed regime.] \citep{esrb_2020}
% [FACTS — Mar-2020 (ESRB 2020), cite exactly:
%   - IM at the four largest EU/UK CCPs rose ~ one third, mainly from CLIENT (not
%     house) portfolios.
%   - daily variation margin on euro-area fund derivative exposures QUINTUPLED;
%     ~6% of funds lacked the pre-stress liquidity to cover cumulative VM.
%   - CMs can apply counterparty-specific IM add-ons up to 50% at short notice;
%     intraday VM is held overnight (a timing asymmetry that delays cash back to
%     payers).]

%==============================================================================
% BEAT 3 — THE GAP: the node that bears this liquidity transformation is the
%               CLIENT-CLEARING TIER (clients -> clearing members -> CCP), and it
%               is the under-studied one. This is where the thesis lives.
% STORY: zoom from "the CCP" to "the tier beneath it" — most work is CCP-centric;
%        the clearing-member / client layer is comparatively dark.
%==============================================================================
\section{The client-clearing tier: the load-bearing, under-studied node}
\label{sec:intro-tier}

% [WRITE: the structural fact] most participants don't face the CCP directly.
[WRITE: Most market participants are not clearing members; they reach the CCP only
THROUGH a member (a bank/FCM or broker) that guarantees their trades, posts
collateral, and absorbs first loss. Clearing is therefore tiered, and the clearing
member is the conduit through which the CCP's margin demands reach end clients —
and through which client distress reaches the CCP.]
\citep{menkveld_vuillemey_2021}

% [WRITE: the gap claim — say it directly] this is what the thesis targets.
[WRITE: State the gap: quantitative CCP research is overwhelmingly CCP-centric and
often static; the clearing-MEMBER / client-clearing perspective is comparatively
under-studied and largely qualitative or regulatory. The economics of the client-
clearing market is explicitly flagged as an open question.] \citep{menkveld_vuillemey_2021}

% [WRITE: three concrete channels that make the tier systemically fragile] each
% gets a citation + a number. These are the "why it matters" beats.
[WRITE: (i) Single-agent dependency — the modal client clears through ONE member
with no backup, so a distressed member can cut clients off from the CCP and force
them to de-risk (Credit Suisse / Archegos).] \citep{ofr_2026}
% [FACTS (OFR 2026): clients = 73% of margin at the largest US/EU CCPs; modal client
%  = single clearing agent; single-agent clients cut cleared positions in stress.]

[WRITE: (ii) Crowding — clearing members' client books are positively correlated,
so correlated portfolios default together in ways member-by-member margin
overlooks; crowding intensifies with volatility.] \citep{huang_menkveld_yu_2021,menkveld_2017}
% [FACTS (Huang, Menkveld & Yu 2021): crowding ~17% of the worst CCP-exposure
%  spikes; top-5 members' share of exposure rises 28% -> 42% in the top-1% tail.]

[WRITE: (iii) Tiering itself is a double-edged structure — it reduces the CCP's
total exposure but raises its expected exposure to the average general clearing
member, i.e. it trades diffuse risk for a concentrated single-point dependency.]
\citep{galbiati_soramaki_2013}
% [Near-miss evidence to deploy here or in 1.2: Nasdaq Clearing 2018 (Aas) tapped
%  EUR 107M of a EUR 166M default fund; ABN AMRO Clearing ~ USD 200M single-trade
%  loss in Mar-2020. Sources flagged VERIFY in the .bib.]

%==============================================================================
% BEAT 4 — THE MECHANISM / THE NOVEL ANGLE: in stress, procyclical margin calls
%               cascade DOWN the tier into client PAYMENT DELAYS; these do not
%               arrive independently — they CLUSTER / self-excite -> a Hawkes lens.
% STORY: this is the thesis's distinctive contribution to the framing. Move from
%        "margin is procyclical" (known) to "the resulting payment failures are
%        self-exciting" (the angle), bridged by the empirical clustering result.
%==============================================================================
\section{Procyclical margin calls, client payment delays, and self-excitation}
\label{sec:intro-paymentdelays}

%  >>> HONESTY FLAG: per writing.md the Hawkes / payment-delay self-excitation
%      mechanism is the thesis's FRAMING LENS and is NOT (yet) in the built model.
%      Write this beat as the mechanism the RQ foregrounds + a proposed modelling
%      contribution. Do NOT state results from it here. If you do implement it,
%      promote it in 1.7; if not, it belongs in RQ + future work. <<<

% [WRITE: the cascade] connect 1.2 + 1.3.
[WRITE: Combine the previous two beats into the mechanism. When volatility spikes,
procyclical margin calls propagate down the tier: the CCP calls its members, members
relay (and add to) the call to clients. A client that cannot meet a call on time
DELAYS or FAILS payment; its member must fund the shortfall, assume and liquidate
the position, and may itself be pushed toward distress — the contagion travels back
UP the tier.] \citep{cont_2017,brunnermeier_pedersen_2009}

% [WRITE: the empirical bridge to self-excitation] breaches are NOT independent.
[WRITE: The key empirical observation that motivates the modelling choice: margin
breaches / clearing-member losses are not independent draws — the most severe losses
cluster and co-move, i.e. one breach raises the probability of the next.]
\citep{jones_perignon_2013,huang_menkveld_yu_2021}

% [WRITE: name the tool] the self-exciting / Hawkes representation.
[WRITE: A natural representation of events whose occurrence raises the short-run
intensity of further events is a self-exciting (Hawkes) point process; such
processes are the standard apparatus for clustered, mutually-exciting financial
events — jump/return contagion and clustered defaults.]
\citep{hawkes_1971,aitsahalia_2015,errais_2010}
% [Optional supporting cite for Hawkes-in-finance breadth: bacry_2015.]
% [POSITIONING: this lets the thesis treat client payment delays as a self-exciting
%  process on the clearing tier — the distinctive lens. State precisely what you
%  will/​won't do with it so the claim matches the implementation.]

%==============================================================================
% BEAT 5 — THE APPROACH + THE PRIMARY CONTRIBUTION: to study propagation through a
%               tier you need a realistic market AND a regulation-grounded clearing
%               layer, coupled. That coupled, calibrated ABM IS THE CONTRIBUTION.
% STORY: justify the method by elimination (static/CCP-centric models can't see the
%        propagation), then plant the flag: the model is the deliverable.
%==============================================================================
\section{An agent-based approach, and why the model is the contribution}
\label{sec:intro-approach}

% [WRITE: why ABM, by elimination] what closed-form / CCP-centric models miss.
[WRITE: Argue the method. The propagation in question is endogenous and path-
dependent: margin calls trigger liquidations, liquidations move prices, price moves
trigger the next margin call. Reduced-form and CCP-centric models that ASSUME the
price process cannot represent this feedback; an agent-based model with an explicit
order book can.] \citep{bookstaber_2014,vytelingum_2025}

% [WRITE: what most CCP ABMs concede, and what this one fixes] the calibration gap.
[WRITE: State the differentiator: most CCP agent-based models assume the price
process; here the market-microstructure layer is CALIBRATED to the empirical
stylised facts of the asset, so the environment the clearing layer acts in is
earned, not assumed.] \citep{gao_2022_xgb,majewski_2018,cont_2001}

% [WRITE: plant the flag — the model itself is the primary contribution] one
% paragraph describing the two-tier architecture at a high level.
[WRITE: Describe the artefact at a glance: a two-tier ABM coupling (i) a calibrated
1-minute call-auction limit-order-book market on E-mini S\&P 500 futures — fundamental,
momentum and zero-intelligence traders around a data-derived (Kalman-filtered)
fundamental — with (ii) a regulation-grounded central-clearing tier: a CCP, banking
and non-banking clearing members, roughly 90 cleared clients, procyclical VaR/SPAN initial
margin, a cover-2 default fund, a five-level loss waterfall, and Almgren--Chriss
fire-sales. The model itself is the central contribution.]
\citep{cont_stoikov_talreja_2008,farmer_2005,almgren_chriss_2000,simudyne_odd,deloitte_ccp}
% [Calibration method, name only here; details in Ch.4: surrogate-assisted simulated
%  method of moments with an XGBoost surrogate, cross-checked by grid search.]
% \citep{gao_2022_hfabm,gao_2023,lamperti_2018,franke_westerhoff_2012}

%==============================================================================
% BEAT 6 — THE RESEARCH QUESTION: state it crisply now that the reader has the
%               tension, the gap, the mechanism, and the tool.
% STORY: everything above converges to one question (+ sub-questions).
%==============================================================================
\section{Research question}
\label{sec:intro-rq}

% [WRITE: the headline RQ — your framing: role of client clearing + tiered
%  procyclical-margin system in propagating/absorbing stress.]
[WRITE: Lead with the main question, e.g.: \emph{What is the role of the tiered
client-clearing structure and its procyclical margin system in propagating or
absorbing market stress through to the CCP?} Phrase as buffer-vs-amplifier: does the
clearing-member layer act as a loss-absorbing buffer between a client and the CCP,
or as a single point of failure that concentrates and transmits stress?]

% [WRITE: sub-questions — keep only what the model actually answers; mark the
%  payment-delay one as the proposed/exploratory strand.]
[WRITE — sub-questions, e.g.:
  (a) Under realistically calibrated stress, is the cover-2 default-fund framework
      sufficient to contain client -> member -> CCP contagion, and what DRIVES a
      breach — the magnitude of the move, or its concentration/speed?
  (b) Does client TYPE (volatility/HFT-like vs directional) or client-book
      concentration change the answer?
  (c) [Exploratory/proposed] If client payment delays are modelled as a self-
      exciting process on the tier, does self-excitation materially change the
      breach point versus i.i.d. payment behaviour?]
% [HONESTY: (a)-(b) are supported by the built model; (c) is the Hawkes extension —
%  keep it as a question/contribution, not a reported result, unless implemented.]

%==============================================================================
% BEAT 7 — CONTRIBUTIONS: enumerate, with the MODEL first.
% STORY: convert the above into a clean, examiner-friendly numbered list.
%==============================================================================
\section{Contributions}
\label{sec:intro-contributions}

% [WRITE: a short lead-in sentence, then the list. Model first, per your framing.]
\begin{enumerate}
  \item % PRIMARY — the artefact.
  [WRITE: A calibrated two-tier ABM that couples an ES-calibrated market-microstructure
  layer to a regulation-grounded clearing tier — addressing the noted gap that little
  quantitative work takes the clearing-member perspective and that most CCP ABMs assume
  rather than calibrate the price process.] \citep{menkveld_vuillemey_2021,gao_2022_xgb}

  \item % methodology realism.
  [WRITE: A regulation-faithful margin/default stack — procyclical VaR/SPAN initial
  margin, a cash/IM capital ratio (CFTC Reg 1.17), a cover-2 Stress-Loss-Over-IM
  default fund (Euronext A9), a deficit-based five-level waterfall, and Almgren--Chriss
  fire-sales — rather than placeholder haircuts.]
  \citep{cftc_reg117,euronext_a9,emir_rts,basel_frtb,almgren_chriss_2000}

  \item % findings on drivers of contagion.
  [WRITE: Findings on the DRIVERS of contagion: cover-2 containment at true COVID
  severity with mutualisation onset only at roughly 2.5$\times$; and the clean structural result that
  contagion tracks the concentration/speed of a move, not its total magnitude.]
  % [FACTS (writing.md; single-seed traces, interim): at actual COVID severity the
  %  cover-2 framework CONTAINS the book (defaulter's own margin/DF absorb it; no
  %  cross-member mutualisation; CCP solvent). Mutualisation (waterfall L3-L4) onsets
  %  at ~2.5x COVID (~ -41%). Gaps-vs-shock: at EQUAL drawdown a single concentrated
  %  intraday shock causes ~3x the client defaults (~29 vs ~9) and reaches CM default +
  %  L4 at c=1 (-19%) vs c~2.5-3 for the overnight-gapped path -> the overnight-gap
  %  structure is a MITIGANT (VM collects between jumps). Label as single-seed/interim;
  %  Monte-Carlo ensemble = future work.]

  \item % methodological contribution.
  [WRITE: A methodological result — a parsimony / robustness finding that richer agent
  dynamics do not robustly improve the fit (over-parameterisation adds estimation
  variance, not fit), and the cautionary lesson that a high surrogate $R^2$ does not
  guarantee a reproducible optimum (the E5 non-replication).]
  % [FACTS: E5's apparent -34% stressed win at screening (D 13.8 -> 9.2) did NOT
  %  replicate at higher resolution (D -> 16.9). Present as a strength: rigorous
  %  robustness study.]

  \item % the angle / proposed extension — keep honest.
  [WRITE — KEEP ONLY IF IMPLEMENTED, else move to future work: A self-exciting
  (Hawkes) representation of client payment delays on the clearing tier, linking the
  empirical clustering of margin breaches to a tractable point-process mechanism.]
  \citep{hawkes_1971,aitsahalia_2015,errais_2010,jones_perignon_2013}
\end{enumerate}

% [Optional methodological caveat to surface honestly somewhere: the stressed
%  market-layer calibration is NOT clearing-invariant (calibrating with the clearing
%  tier active alters the stressed fit, D 13.8 -> 28) — a finding in its own right.]

%==============================================================================
% BEAT 8 — ROADMAP: one short paragraph mapping the chapters.
% STORY: tell the reader the route. Keep to prose, not a list, if your style guide
%        prefers (writing.md leans prose).
%==============================================================================
\section{Thesis structure}
\label{sec:intro-roadmap}

% [WRITE: a single guiding paragraph. Suggested mapping (adjust to your chapters):]
[WRITE: Chapter~\ref{ch:litreview} reviews the three literatures the thesis sits at
the junction of — agent-based market microstructure, CCP / clearing risk and
contagion, and ABM calibration. Chapter~\ref{ch:model} specifies the two-tier model
(the ODD protocol, agents, the call-auction LOB, the fundamental, and the clearing
tier). Chapter~\ref{ch:data} covers the ES data, the calm/stressed regimes, and the
calibration method. Chapter~\ref{ch:market-results} validates the market layer against
the empirical stylised facts and reports the robustness campaign. Chapter~\ref{ch:clearing-results}
runs the clearing-layer and contagion experiments. Chapter~\ref{ch:discussion}
discusses drivers, methodological findings, and limitations, and Chapter~\ref{ch:conclusion}
answers the research question and sets out future work.]
% [Update the \ref labels to match your chapter files. Placeholders only.]
